% !PRINTED 32
\weberpage{32}

\begin{center}
{\large \S.\ 5.}

\vspace{0.5em}
{\large\gesperrt{Fractional Functions; Divisibility.}}
\end{center}

If $F(x)$ and $f(x)$ are two whole rational functions of $x$, then
the fraction
\begin{equation}
 {F(x)\over f(x)}
\tag{1}
\end{equation}
is called a \gesperrt{fractional rational}, or briefly a
\gesperrt{fractional} or \gesperrt{rational function} of $x$.

If the degree of the numerator is lower than the degree of the
denominator, the function is called a \gesperrt{proper fraction}; in
the opposite case it is called an \gesperrt{improper fraction}.

According to \S~3, the whole rational functions $Q$ and $\varphi(x)$
can be determined so that
\begin{equation}
 F(x)=Qf(x)+\varphi(x),
\tag{2}
\end{equation}
and the degree of $\varphi(x)$ is smaller than the degree of $f(x)$;
consequently
\begin{equation}
 {F(x)\over f(x)}=Q+{\varphi(x)\over f(x)}
\tag{3}
\end{equation}
and hence the theorem:

\gesperrt{Every fractional function can be decomposed into the sum of
a whole function and a proper fractional function.}

If $n$ is the degree of $f(x)$, then the degree of $\varphi(x)$ is at
most $n-1$; it may also be lower. In particular, the case may occur in
which $\varphi(x)$ \gesperrt{vanishes identically}.

\gesperrt{In this case the function $F(x)$ is said to be divisible by
$f(x)$.} The function
\[
 {F(x)\over f(x)}
\]
is in this case only apparently fractional; in reality it is equal to
the whole function $Q$.

The formula
\[
 {x^m-1\over x-1}=x^{m-1}+x^{m-2}+x^{m-3}+\cdots+1
\]
gives a simple example.
